% Starter deck -- edit me, then `make` (-> talk.pdf). Lives in the extension so a
% vendored engine (the framework) packs it into any livery's Overleaf zip (overleaf_zip.py).
% Every feature of the theme is demonstrated in examples/template.tex.
\documentclass[aspectratio=169]{beamer}

% livery=<dir>|brand=ncar|ucar|ucp|ncarucar   title=dark|light   footer=minimal|full   mathfont=serif|stix|sans|...
\usetheme[brand=ncar, title=dark]{Engine}

\usepackage{booktabs}
\usepackage{listings}

\title{Talk Title Goes Here}
\subtitle{An optional subtitle}
\author{Your Name}
\institute{Your Title, Your Lab, Your Institution}
\date{\today}
\titlegraphic{\includegraphics{images/wallpaper/mesa_lab_sunset.jpg}} % optional

\begin{document}

\begin{frame}[plain,noframenumbering]
  \titlepage
\end{frame}

\section{Introduction}

\begin{frame}{A Content Slide}
  \begin{itemize}
  \item Make a point
    \begin{itemize}
    \item and support it
    \end{itemize}
  \item Highlight \alert{one key idea}
  \end{itemize}
  \begin{block}{A block}
    Blocks, example blocks, and alert blocks are all branded.
  \end{block}
\end{frame}

\begin{frame}[enginebg=images/wallpaper/mesa_lab.jpg]{A Photo Slide}
  \begin{enginefeaturepanel}
    \begin{block}{The body on a panel}
      \texttt{\textbackslash begin\{frame\}[enginebg=photo.jpg]}, the body in \texttt{enginefeaturepanel}
    \end{block}
  \end{enginefeaturepanel}
\end{frame}

\section{Details}

\begin{frame}[fragile]{Code}
\begin{lstlisting}[language=bash]
#!/bin/bash
# the whole deck is one file; edit, then `make`
latexmk -xelatex talk.tex
\end{lstlisting}
\end{frame}

\engineclosingframe{Thank you!}{%
  Questions?\par\medskip
  \textbf{Your Name}\\ \texttt{you@example.org}}

\end{document}
